\subsection{Zeichnung und Berechnung einer Kollektorschaltung}
% Alt: Praktikumsaufgabe 9
\Aufgabe{
    \label{Aufg36}
    \begin{itemize}
        \item[a)] Zeichnen Sie die folgende Schaltung:
              Beschalten Sie einen npn-Transistor mit einem Basisvorwiderstand von $R_\mathrm{B} = 4,7\,\mathrm{k\Omega}$, einem 
              Emitterwiderstand von  $R_\mathrm{E} = 1\,\mathrm{k\Omega}$ und einer Spannungsversorgung von $U_\mathrm{0} = 15\,\mathrm{V}$.
              Schalten Sie eine $10 \,\mathrm{V}$-Z-Diode in Sperrrichtung von der Basis auf Masse. \newline
              Verwenden Sie nur eine Spannungsversorgung, welche sowohl die Basis als auch den Kollektor versorgt.
        \item[b)] Berechnen Sie die Emitterspannung $U_\mathrm{E}$, und den Emitterstrom $I_\mathrm{E}$.
        \item[c)] Berechnen Sie die Leistung am Emitterwiderstand $P_\mathrm{RE}$.
    \end{itemize}
}


\Loesung{
    \begin{itemize}
        \item[a)] Schaltung:
              \begin{figure}[H]
                  \centering
                  \input{Tikz/src/LsgPraktikumZeichnung.tex}
                  \alttext{Die Abbildung zeigt eine Transistorschaltung.
                  Die Versorgungsspannung wird durch die Spannungsquelle \(U_0\) bereitgestellt, die links zwischen dem oberen und dem unteren Versorgungsknoten eingezeichnet ist.
                  Rechts befindet sich der NPN-Transistor. Vom Emitter des Transistors führt der Widerstand \(R_\mathrm{E}\) zum unteren Versorgungsknoten.
                  Die Basis des Transistors ist mit einem Knoten verbunden, der über den Widerstand \(R_\mathrm{B}\) an den oberen Versorgungsknoten angeschlossen ist. Zwischen diesem Basisknoten und dem unteren Versorgungsknoten befindet sich eine Z-Diode. Die Z-Diode liegt somit parallel zur Basis-Emitter-Strecke.
                  Der obere Versorgungsknoten ist sowohl mit dem Kollektor des Transistors als auch mit dem Widerstand \(R_\mathrm{B}\)verbunden. Der untere Versorgungsknoten verbindet den Emitterwiderstand, die Z-Diode und Spannungsquelle.}
                  \label{fig:LsgPraktikumZeichnung}
              \end{figure}
              
              
        \item[b)] Berechnung der Emitterspannung $U_\mathrm{E}$ \\
              Die Zenerdiode fixiert die Basisspannung des Transistors auf $U_\mathrm{Z} = 10\,\mathrm{V}$. 
              Die Basis-Emitter-Spannung beträgt $U_{BE} = 0,7\,\mathrm{V}$.
              \begin{gather*}
                  U_\mathrm{E} = U_\mathrm{Z} - U_\mathrm{BE} = 10\,\mathrm{V} - 0,7\,\mathrm{V} = 9,3\,\mathrm{V}                  \\
                  I_\mathrm{E} = \frac{U_\mathrm{E}}{R_\mathrm{E}} = \frac{9,3\,\mathrm{V}}{1\,\mathrm{k\Omega}} = 9,3\,\mathrm{mA}
              \end{gather*}
              
              
        \item[c)] Berechnung der Leistung am Emitterwiderstand $P_{RE}$ \\
              $$ P_\mathrm{RE} = I_\mathrm{E}^2 \cdot R_\mathrm{E} = (9,3\,\mathrm{mA})^2 \cdot 1\,\mathrm{k\Omega} = 86,49\,\mathrm{mW} $$
              
    \end{itemize}
    Siehe: Abschnitt \ref{fig:GrundschaltungenBipolartransistor}
}